\section{Codex of Potential (Meta Progression)}

\subsection{Purpose and Scope}

The \emph{Codex of Potential} defines the primary account-wide meta-progression system.
It does not directly increase hero power, but instead unlocks access to systems, mechanics,
tiers, and options across the game.

The Codex is driven exclusively by deliberate player actions and does not rely on time-based
progression, currency investment, or passive accumulation.

Its core responsibility is to gate and sequence the player's long-term possibilities
while remaining orthogonal to Hero progression, Orbit progression, and moment-to-moment combat balance.

\subsubsection{System Overview}

\begin{center}
	\begin{tabular}{|l|l|l|l|}
		\hline
		\textbf{Input Source} & \textbf{Tracked Signal} & \textbf{Codex Element} & \textbf{Result} \\
		\hline
		Combat System & Enemy kills (rank / tags) & Deed Requirements & Progress increment \\
		\hline
		Crafting System & Crafts / Refines & Deed Requirements & Progress increment \\
		\hline
		Map System & Map and Wave completions & Deed Requirements & Progress increment \\
		\hline
		Difficulty System & Difficulty thresholds & Deed Requirements & Completion gating \\
		\hline
		Codex State & Completed Deeds & Unlock Registry & Global unlocks \\
		\hline
	\end{tabular}
\end{center}

The Codex of Potential acts purely as a progression controller.
It consumes gameplay signals but never injects combat effects directly.

\subsection{Structural Model}

The Codex is organized into multiple \emph{Chapters} (internally referred to as \emph{Pillars}),
each representing a fundamental domain of progression (e.g. Crafting, Heroes, World Tiers).

Each Chapter contains a predefined set of \emph{Deeds}, arranged in a structured and deterministic layout.
Dependencies between Deeds are explicit and resolved through directed parent relationships.

The overall structure is static and authored; Deeds are never randomized or reshuffled.

\subsubsection{Chapter Examples}

The following tables provide concrete examples of Chapters and their associated Deeds.
They are illustrative and not exhaustive, but represent the intended progression flow
and unlock pacing within each Chapter.

\paragraph{Crafting Chapter}

\begin{center}
	\begin{tabular}{p{3.2cm}p{4cm}p{8cm}}
		\textbf{Deed} & \textbf{Short Description} & \textbf{Unlocks / Effects} \\
		\hline
		Basic Gear &
		\emph{``Let's see what we can do with these shreds.''} &
		Unlocks basic gear usage and simple item crafting from common loot. Establishes the baseline of the item system. \\
		\hline
		Infuse I &
		\emph{``We can make use of these magical remains.''} &
		Unlocks the first tier of infusion: adding simple magical properties to base items using low-tier materials. \\
		\hline
		Re-Infuse &
		Allows re-rolling or replacing existing infusions at a cost. &
		Unlocks re-infusing items, enabling early deterministic rerolls instead of pure item hunting. \\
		\hline
		Gear II &
		\emph{``Find some of the rarer materials, so we can improve the gear.''} &
		Unlocks access to higher gear tiers and rarer crafting materials; raises the ceiling for item power. \\
		\hline
		Refine / Tempering &
		Introduces finer control over item stats via tempering. &
		Unlocks refining and tempering operations that adjust existing stats within defined ranges. \\
		\hline
		Infuse II &
		Second tier of infusion for more powerful effects or additional affix slots. &
		Unlocks stronger infusions and/or more slots per item, increasing build complexity. \\
		\hline
		Gear III &
		Unlocks late-game gear tiers. &
		Gives access to high-end gear bases and top-tier materials for endgame items. \\
		\hline
		Infuse III &
		Third tier of infusion, oriented towards endgame builds. &
		Unlocks the highest tier of infusions with strong, build-defining properties. \\
		\hline
		Over-Tempering &
		Extreme, high-risk / high-reward tempering. &
		Unlocks over-tempering mechanics that allow pushing items beyond normal limits at significant cost or risk. \\
	\end{tabular}
\end{center}


\paragraph{World and Map Chapter}

\begin{center}
	\begin{tabular}{p{3.2cm}p{4cm}p{8cm}}
		\textbf{Deed} & \textbf{Short Description} & \textbf{Unlocks / Effects} \\
		\hline
		Empower Maps &
		Enables basic Map empowerment. &
		Unlocks modifiers that increase Map difficulty in exchange for better rewards (e.g. more enemies, higher damage or HP, better loot). \\
		\hline
		Map Live-Events &
		Introduces dynamic events during Maps. &
		Unlocks live events during a Map (e.g. reinforcements, shrines, hazards) that alter Wave flow and risk/reward mid-run. \\
		\hline
		Undead Monsters &
		Expands the enemy ecosystem with undead. &
		Unlocks undead-themed enemy packs, bosses and tags; adds new resistances, weaknesses, and synergies. \\
		\hline
		Holy Monsters &
		Introduces holy or fanatic enemy types. &
		Unlocks holy-themed enemies and factions with their own mechanics and counters, further diversifying encounters. \\
	\end{tabular}
\end{center}


\paragraph{Hero Chapter}

\begin{center}
	\begin{tabular}{p{3.2cm}p{4cm}p{8cm}}
		\textbf{Deed} & \textbf{Short Description} & \textbf{Unlocks / Effects} \\
		\hline
		Hero Packs (Story / Theme) &
		Unlocks groups of Heroes along narrative or thematic lines. &
		Frees 2--3 Heroes at a time, grouped by story or theme, instead of single isolated Heroes. Encourages experimenting with new team archetypes. \\
		\hline
		Skill Tiers &
		Unlocks deeper layers of the skill trees. &
		Opens higher skill tiers (e.g. Tier 2/3 nodes) for existing Heroes, enabling more specialised builds and advanced synergies. \\
		\hline
		Crown Sockets &
		Prepares Heroes for the endgame Crown system. &
		Unlocks Crown sockets on Heroes or their skill trees, allowing equipping powerful endgame Crown boosts once available. \\
		\hline
		Orbit Access &
		Grants access to the Orbit progression interface. &
		Unlocks the Orbit system for Heroes; further Orbit progression unfolds inherently through Hero advancement. \\
	\end{tabular}
\end{center}


\subsection{Deeds and Active Focus}

A \emph{Deed} represents a deliberate, meaningful action that can be completed through normal gameplay.
Deeds define explicit completion requirements, such as kill counts, map completions, or difficulty thresholds.

At any given time, a limited number of Deeds can be assigned as the player's \emph{Active Focus}.
Only Deeds in Active Focus receive progress updates from gameplay events.

Once a Deed is completed, it occupies its Active Focus slot until its reward is consciously claimed,
at which point the slot becomes available for a new Deed.

\paragraph{Initial Deed (Onboarding)}

To introduce the Codex mechanics, the first Deed is automatically assigned as an Active Focus.
This initial Deed serves purely as an onboarding element, demonstrating the Codex workflow
(selection, progress, completion, and claiming) once before full player control is granted.


\subsubsection{Deed Structure}

\begin{center}
	\begin{tabular}{|l|l|l|l|}
		\hline
		\textbf{Deed Type} & \textbf{Requirement Signal} & \textbf{Aggregation} & \textbf{Typical Unlock} \\
		\hline
		Combat Deed & Enemy kills & Weighted by rank & Crafting feature access \\
		\hline
		Exploration Deed & Map completions & Per difficulty & World tier unlock \\
		\hline
		Crafting Deed & Craftings \& reifnements & per Tier (output) & Crafting tier \& recipes access \\
		\hline
		Challenge Deed & Elite / Boss kills & Absolute count & System feature unlock \\
		\hline
		Progression Deed & Mixed signals & AND-combined & Hero or skill access \\
		\hline
	\end{tabular}
\end{center}

Deeds may combine multiple requirement signals.
All requirements must be fulfilled for completion.

\subsubsection{Active Focus Slots}

The Codex supports a limited number of \emph{Active Focus Slots}.
Each slot may hold exactly one Deed.

\begin{itemize}
	\item The player starts with a single Active Focus Slot.
	\item Additional slots are unlocked through Codex progression.
	\item Each slot tracks progress independently.
\end{itemize}
Multiple Deeds may be completed in parallel if sufficient slots are available.

\subsubsection{Repeatable Deeds}
Certain deeds may be marked as repeatable to support endgame progression.
Repeatable deeds:
\begin{itemize}
	\item Reset their progress after being claimed
	\item Continue to grant rewards on subsequent completions
\end{itemize}
Repeatability is balanced through scaling requirements rather than reward reduction.
Each repetition increases requirement thresholds slightly (e.g. +5\% per completion), ensuring:
\begin{itemize}
	\item Sustained progression without infinite farming
	\item No abrupt or unintuitive reward drops
\end{itemize}
This approach preserves reward consistency while naturally increasing challenge over time.

\subsection{Progress Tracking}

Progress towards a Deed is tracked incrementally and persists independently of whether the Deed
remains actively focused or is temporarily inactive.

Progress evaluation is requirement-based and supports aggregated conditions
(e.g. weighted kills, tagged enemies, map completions by difficulty).

Completion checks are deterministic and occur immediately when requirements are met.

\subsubsection{Progress Evaluation}

\paragraph{Progress Persistence}

Progress towards a Deed is stored persistently on a per-Deed basis.

Switching the Active Focus:
\begin{itemize}
	\item does not reset progress,
	\item does not invalidate partial completion,
	\item does not trigger completion checks while inactive.
\end{itemize}

This allows players to pause and resume Deeds freely without penalty.

\subsection{Unlock Propagation}

Upon completion, a Deed emits one or more unlock events.
Unlocks are declarative and are registered centrally in the Unlock Registry.

Other systems query unlock states directly and do not maintain duplicated progression logic.
The Codex itself does not apply gameplay effects beyond signaling completed unlocks.

\paragraph{Unified Unlock Handling}
The Codex Unlock system is designed as a generic unlock pipeline and is not limited to deed progression.
Unlocks are processed in two stages:

\begin{itemize}
	\item The Unlock Registry records that an unlock condition has been fulfilled (logical state).
	\item Unlock Consumers apply gameplay effects based on the unlock type (system-side effects).
\end{itemize}
This mechanism can be reused for unlocks outside of deeds, such as:
\begin{itemize}
	\item Unlocking crafting tiers
	\item Unlocking skill recipes
	\item Unlocking system features or UI elements
\end{itemize}
Importantly, runtime conditions (e.g. hero level) act as gates and are not treated as unlock events themselves. Unlocks define \emph{what} becomes available; runtime conditions define \emph{when} it can be used.

\subsubsection{Completion and Claiming}

Completion and reward claiming are intentionally separated.

\begin{itemize}
	\item A Deed is considered \emph{completed} once all requirements are met.
	\item Completion immediately registers unlocks in the Unlock Registry.
	\item Claiming the Deed reward is a conscious player action.
\end{itemize}
A completed Deed continues to occupy its Active Focus Slot
until its reward is claimed.

\paragraph{Coupon Rewards via Difficulty Score}
Deeds grant a single reward currency: Coupons.
The amount of Coupons is not configured per deed but derived from a centralized reward calculation.
Each deed has an implicit Difficulty Score, calculated from its requirements using:
\begin{itemize}
	\item Requirement-specific weights (enemy rank, item tier, map difficulty)
	\item Category multipliers (e.g. kills, maps, crafts, modifies)
\end{itemize}
Coupons are awarded on claim using a global conversion factor:
\[
\text{Coupons} = \frac{\text{Difficulty Score}}{30}
\] \\
The final output is normalized using a rounding ladder to produce clean, readable reward values (e.g. 25, 50, 100, 250, 500, 1000).
An optional global cap (e.g. 2500 Coupons) prevents excessive rewards at the upper end.\\
All reward tuning is handled centrally, allowing global balance changes without modifying individual deeds.

\subsection{Non-Goals}

The Codex of Potential explicitly does not:
\begin{itemize}
	\item provide direct numerical power increases,
	\item replace Hero progression or Orbit-based specialization,
	\item function as a passive or time-gated system,
	\item enforce mandatory progression to access basic gameplay loops.
\end{itemize}
